\documentclass[11pt, a4paper]{article}

% Gemeinsame Style-Definitionen (TEXINPUTS muss templates/ enthalten — siehe scripts/build_tex.py)
% bfw-style.tex — gemeinsame Style-Definitionen für alle Handouts/Aufgabenhefte
% Voraussetzung: Kompilation mit xelatex (wegen fontspec)

\usepackage[ngerman]{babel}
\usepackage{geometry}
\geometry{a4paper, top=2.2cm, bottom=2.2cm, left=2.3cm, right=2.3cm}

\usepackage{fontspec}
% Fallback-Kette: Source Sans Pro -> Calibri -> Segoe UI -> Latin Modern Sans
% (Latin Modern ist immer mit MiKTeX vorhanden)
\IfFontExistsTF{Source Sans Pro}{%
  \setmainfont{Source Sans Pro}[Ligatures=TeX]%
}{\IfFontExistsTF{Calibri}{%
  \setmainfont{Calibri}[Ligatures=TeX]%
}{\IfFontExistsTF{Segoe UI}{%
  \setmainfont{Segoe UI}[Ligatures=TeX]%
}{%
  \setmainfont{Latin Modern Sans}[Ligatures=TeX]%
}}}
\IfFontExistsTF{Source Code Pro}{%
  \setmonofont{Source Code Pro}[Scale=0.92]%
}{\IfFontExistsTF{Consolas}{%
  \setmonofont{Consolas}[Scale=0.92]%
}{%
  \setmonofont{Latin Modern Mono}[Scale=0.92]%
}}

\usepackage{xcolor}
% Farbpalette: hell, freundlich, modern
\definecolor{bfwPrimary}{HTML}{0EA5A4}    % Petrol/Türkis - Primär
\definecolor{bfwPrimaryDark}{HTML}{0F766E} % dunkler Petrol für Headings
\definecolor{bfwAccent}{HTML}{F59E0B}     % Amber - Akzent
\definecolor{bfwSuccess}{HTML}{10B981}    % Grün - Erfolg / Lösung
\definecolor{bfwWarn}{HTML}{F97316}       % Orange - Achtung
\definecolor{bfwInk}{HTML}{0F172A}        % fast Schwarz
\definecolor{bfwInkSoft}{HTML}{475569}    % gedämpftes Grau
\definecolor{bfwBgSoft}{HTML}{F8FAFC}     % off-white
\definecolor{bfwBgTeal}{HTML}{ECFEFF}     % helles Türkis
\definecolor{bfwBgAmber}{HTML}{FEF3C7}    % helles Gelb
\definecolor{bfwBgRose}{HTML}{FFE4E6}     % helles Rosé für Eselsbrücken

\usepackage[colorlinks=true, linkcolor=bfwPrimaryDark, urlcolor=bfwPrimaryDark]{hyperref}
\usepackage{enumitem}
\setlist[itemize]{leftmargin=*, itemsep=2pt, topsep=2pt}
\setlist[enumerate]{leftmargin=*, itemsep=2pt, topsep=2pt}

\usepackage[most]{tcolorbox}
\usepackage{booktabs}
\usepackage{tikz}
\usetikzlibrary{decorations.pathmorphing, positioning, shapes.geometric, arrows.meta}

% Merkkästchen — wichtige Definition / Kernpunkt
\newtcolorbox{merksatz}[1][]{%
  enhanced, breakable,
  colback=bfwBgTeal,
  colframe=bfwPrimary,
  boxrule=0.6pt,
  arc=4pt,
  left=10pt, right=10pt, top=8pt, bottom=8pt,
  fonttitle=\bfseries\color{bfwPrimaryDark},
  title={\faLightbulb\ Merksatz},
  #1
}

% Eselsbrücke — Mnemoniks
\newtcolorbox{eselsbruecke}[1][]{%
  enhanced, breakable,
  colback=bfwBgRose,
  colframe=bfwAccent,
  boxrule=0.6pt,
  arc=4pt,
  left=10pt, right=10pt, top=8pt, bottom=8pt,
  fonttitle=\bfseries\color{bfwWarn},
  title={Eselsbrücke},
  #1
}

% Beispiel-Box
\newtcolorbox{beispiel}[1][]{%
  enhanced, breakable,
  colback=bfwBgSoft,
  colframe=bfwInkSoft,
  boxrule=0.4pt,
  arc=3pt,
  left=10pt, right=10pt, top=8pt, bottom=8pt,
  fonttitle=\bfseries\color{bfwInk},
  title={Beispiel},
  #1
}

% Lösungs-Box (für Aufgabenhefte)
\newtcolorbox{loesung}[1][]{%
  enhanced, breakable,
  colback=bfwBgAmber!40,
  colframe=bfwSuccess,
  boxrule=0.6pt,
  arc=3pt,
  left=10pt, right=10pt, top=8pt, bottom=8pt,
  fonttitle=\bfseries\color{bfwSuccess!70!black},
  title={Lösung},
  #1
}

% Glossar-Eintrag
\newcommand{\glossareintrag}[2]{%
  \par\noindent\textbf{\color{bfwPrimaryDark}#1} \enspace #2\par\smallskip
}

% Section-Stil — etwas farbiger als Standard
\usepackage{titlesec}
\titleformat{\section}
  {\Large\bfseries\color{bfwPrimaryDark}}
  {\thesection}{0.6em}{}
\titleformat{\subsection}
  {\large\bfseries\color{bfwPrimary}}
  {\thesubsection}{0.5em}{}
\titleformat{\subsubsection}
  {\normalsize\bfseries\color{bfwInk}}
  {\thesubsubsection}{0.4em}{}

% TikZ-Stil "skizzenhaft" — etwas weicher als Rough.js, aber stabil
\tikzset{
  roughbox/.style = {
    draw=bfwPrimaryDark, thick, fill=bfwBgTeal,
    rounded corners=4pt, inner sep=6pt
  },
  roughaccent/.style = {
    draw=bfwAccent, thick, fill=bfwBgAmber,
    rounded corners=4pt, inner sep=6pt
  },
  roughline/.style = {
    draw=bfwInkSoft, thick, -{Stealth[length=2.5mm]}
  }
}

% Bullet-Symbole (bewusst kein FontAwesome — vermeidet Font-Installations-Probleme)
\newcommand{\faLightbulb}{$\bullet$}
\newcommand{\faBookmark}{$\bullet$}

% Header/Footer
\usepackage{fancyhdr}
\pagestyle{fancy}
\fancyhf{}
\renewcommand{\headrulewidth}{0pt}
\fancyfoot[L]{\small\color{bfwInkSoft}\thepage}
\fancyfoot[R]{\small\color{bfwInkSoft} BFW M-064 Beschaffung im Büromanagement}


\title{\vspace*{-1.5cm}\Huge\bfseries\color{bfwPrimaryDark}Session~2: Der Betrieb in der Wirtschaft\\[0.4em]
  \large\color{bfwInkSoft}\textnormal{Bedürfnisse, Bedarf \& Wirtschaftskreislauf}}
\author{\textsf{Lernfeld 1 --- Die eigene Rolle im Betrieb mitgestalten \textperiodcentered{} \small\copyright\ Can Siebert\,\textperiodcentered\,2026}}
\date{}

\begin{document}
\maketitle
\thispagestyle{empty}

\vspace{1em}
\begin{merksatz}[title={\bfseries Worum geht es in dieser Session?}]
Jedes Wirtschaften beginnt mit einem \emph{Bedürfnis} --- einem Mangelgefühl, das der
Mensch beseitigen möchte. Aus Bedürfnissen wird über den \emph{Bedarf} eine \emph{Nachfrage},
die auf ein knappes Angebot an \emph{Gütern} trifft. Weil Güter knapp sind, muss der Mensch
planvoll wirtschaften --- nach dem \emph{ökonomischen Prinzip}. Dazu werden
\emph{Produktionsfaktoren} kombiniert. Wie all diese Größen zwischen Haushalten,
Unternehmen, Banken, Staat und Ausland zusammenspielen, zeigt der \emph{Wirtschaftskreislauf}.
Diese Grundbegriffe der Volkswirtschaft sind die Basis für das gesamte Lernfeld und
für die IHK-Abschlussprüfung.
\end{merksatz}

\tableofcontents
\vspace{1em}
\hrule
\vspace{1em}

\section{Bedürfnisse --- der Ausgangspunkt allen Wirtschaftens}

Ein \textbf{Bedürfnis} ist das Gefühl eines Mangels, verbunden mit dem Wunsch, diesen
Mangel zu beseitigen. Bedürfnisse lassen sich nach mehreren Kriterien einteilen.

\subsection{Einteilung nach der Dringlichkeit}

Nach der \emph{Dringlichkeit der Bedürfnisbefriedigung} unterscheidet man drei Gruppen:

\begin{itemize}
  \item \textbf{Existenzbedürfnisse} (Grundbedürfnisse) --- sie sichern das Überleben,
    z.\,B.\ Essen, Trinken, Kleidung und Wohnung.
  \item \textbf{Kulturbedürfnisse} --- sie werden von der Umwelt bzw.\ dem Kulturkreis
    geprägt und dort als üblich angesehen, z.\,B.\ Bildung, modische Kleidung, Kino,
    Theater oder ein Hobby.
  \item \textbf{Luxusbedürfnisse} --- nicht notwendige Ansprüche, die sich nicht jeder
    leisten kann, z.\,B.\ eine Yacht, ein Sportwagen oder eine Villa.
\end{itemize}

Die Grenzen zwischen diesen Bedürfnissen sind \emph{fließend}; eine konkrete Abgrenzung
ist nicht immer möglich (ein Smartphone kann für den einen Kulturbedürfnis, für den
anderen Luxus sein).

\subsection{Weitere Einteilungen der Bedürfnisse}

Neben der Dringlichkeit lassen sich Bedürfnisse nach weiteren Kriterien ordnen:

\begin{center}
\begin{tabular}{p{4.5cm} p{9cm}}
\toprule
\textbf{Kriterium} & \textbf{Unterscheidung}\\
\midrule
Gegenstand & \emph{materielle} Bedürfnisse (Smartphone, Schuhe) vs.\ \emph{immaterielle} Bedürfnisse (Anerkennung, eine Nachricht schreiben)\\
Grad der Bewusstheit & \emph{akute} Bedürfnisse (offener Wunsch) vs.\ \emph{latente} Bedürfnisse (schlummernd, z.\,B.\ durch Werbung geweckt)\\
Möglichkeit der Befriedigung & \emph{Individualbedürfnis} (geht von einem Menschen aus) vs.\ \emph{Kollektivbedürfnis} (entsteht aus dem Zusammenleben, z.\,B.\ Umweltschutz, gutes Autobahnnetz)\\
\bottomrule
\end{tabular}
\end{center}

\begin{beispiel}
Ein „schlummernder" Wunsch nach neuen Schuhen (\emph{latent}) kann zu einem \emph{akuten}
Bedürfnis werden, wenn Sie in einem Onlineshop ein Paar entdecken, das Sie unbedingt
besitzen möchten.
\end{beispiel}

\subsection{Die Bedürfnispyramide nach Maslow}

Der Psychologe Abraham \textbf{Maslow} ordnete die menschlichen Bedürfnisse in einer
\emph{Pyramide} an. Die Grundannahme: Erst wenn die unteren, dringlicheren Bedürfnisse
weitgehend befriedigt sind, treten die höheren in den Vordergrund. Von unten nach oben:

\begin{enumerate}
  \item \textbf{Physiologische Grundbedürfnisse} --- Essen, Trinken, Schlaf.
  \item \textbf{Sicherheitsbedürfnisse} --- Wohnung, Arbeitsplatz, Gesundheit.
  \item \textbf{Soziale Bedürfnisse} --- Freundschaft, Zugehörigkeit, Liebe.
  \item \textbf{Individualbedürfnisse (Wertschätzung)} --- Anerkennung, Status, Erfolg.
  \item \textbf{Selbstverwirklichung} --- Entfaltung der eigenen Fähigkeiten.
\end{enumerate}

Die untersten drei Stufen gelten als \emph{Defizitbedürfnisse} (ihr Fehlen wird als Mangel
empfunden), die oberen als \emph{Wachstumsbedürfnisse}. Das Modell ist anschaulich, aber
vereinfachend --- in der Realität verlaufen die Übergänge fließend.

\section{Vom Bedürfnis über den Bedarf zur Nachfrage}

Nicht jedes Bedürfnis wird zum Kauf. Erst wenn für die Befriedigung eines Bedürfnisses die
nötigen \emph{finanziellen Mittel} (Einkommen bzw.\ Kaufkraft) vorhanden sind, spricht man
von \textbf{Bedarf}. Wird das Gut dann tatsächlich gekauft, entsteht \textbf{Nachfrage}.

\begin{merksatz}
\textbf{Bedürfnis} (Mangelgefühl) $\rightarrow$ \textbf{Bedarf} (mit Kaufkraft ausgestattetes
Bedürfnis) $\rightarrow$ \textbf{Nachfrage} (marktwirksames Bedürfnis, tatsächlicher Kauf).
Die Kaufkraft ist das Bindeglied: Ohne Einkommen bleibt das Bedürfnis wirkungslos.
\end{merksatz}

\begin{eselsbruecke}
\textbf{„Erst wünschen, dann können, dann kaufen."} --- Bedürfnis (wünschen),
Bedarf (können = Kaufkraft), Nachfrage (kaufen).
\end{eselsbruecke}

\subsection{Der Markt: Angebot und Nachfrage}

Auf dem \textbf{Markt} treffen \emph{Angebot} (die zum Verkauf bereitgestellten Güter der
Anbieter) und \emph{Nachfrage} (die kaufkräftigen Wünsche der Nachfrager) zusammen. Aus
ihrem Zusammenspiel bildet sich der \emph{Preis}. Grundidee der Preisbildung:

\begin{itemize}
  \item Ist ein Gut sehr knapp oder die Nachfrage sehr groß, \emph{steigt} der Preis.
  \item Ist ein Gut reichlich vorhanden oder die Nachfrage gering, \emph{sinkt} der Preis.
\end{itemize}

Der Preis hat damit eine \emph{Lenkungsfunktion}: Er signalisiert Knappheit und steuert,
welche Güter in welcher Menge produziert und nachgefragt werden.

\section{Güterarten}

Alle Mittel, mit denen der Mensch seine Bedürfnisse befriedigt, nennt man \textbf{Güter}.
Güter werden nach mehreren Merkmalen geordnet.

\subsection{Freie und knappe Güter}

\textbf{Freie Güter} sind in großen Mengen vorhanden und haben keinen Preis, z.\,B.\ Luft,
Sonne oder Meerwasser. \textbf{Knappe Güter (Wirtschaftsgüter)} stehen nur begrenzt zur
Verfügung, müssen durch wirtschaftliche Tätigkeit erzeugt oder bereitgestellt werden,
verursachen Kosten und haben deshalb einen \emph{Preis}. Je knapper ein Gut bzw.\ je größer
die Nachfrage, desto höher der Preis.

\subsection{Weitere Unterscheidungen knapper Güter}

\begin{itemize}
  \item \textbf{Materielle Güter (Sachgüter)} --- fassbar, z.\,B.\ Schreibtische,
    Lebensmittel. \\
    \textbf{Immaterielle Güter} --- nicht fassbar; sie teilen sich in \emph{Dienstleistungen}
    (Handlungen mit Nutzen, z.\,B.\ Reparaturen, Beratungen) und \emph{Rechte}
    (z.\,B.\ Warenzeichen/Marken, Patente).
  \item \textbf{Konsumgüter} --- von Endverbrauchern (Konsumenten) genutzt, z.\,B.\ das
    private Smartphone. \\
    \textbf{Produktionsgüter} --- von Unternehmen im Herstellungsprozess eingesetzt,
    z.\,B.\ Maschinen, Betriebsfahrzeug, dienstliches Smartphone.
  \item \textbf{Gebrauchsgüter} --- mehrmals verwendbar, z.\,B.\ eine Maschine. \\
    \textbf{Verbrauchsgüter} --- nur einmal verwendbar, z.\,B.\ Holz als Rohstoff.
\end{itemize}

\begin{center}
\begin{tabular}{p{3cm} p{5.5cm} p{5.5cm}}
\toprule
\textbf{Beziehung} & \textbf{Komplementärgüter} & \textbf{Substitutionsgüter}\\
\midrule
Wirkung & ergänzen sich, werden gemeinsam nachgefragt & ersetzen sich, sind austauschbar\\
Beispiel & Maschine \& Strom, Butter \& Toast & Butter \emph{oder} Margarine, Lkw \emph{oder} Bahn\\
Nachfrage & steigt eines, steigt meist auch das andere & verläuft gegenläufig (steigt Butter, sinkt Margarine)\\
\bottomrule
\end{tabular}
\end{center}

\begin{eselsbruecke}
\textbf{Komplementär = Kompagnon} (Partner, die zusammengehören) ---
\textbf{Substitut = Sub} (Ersatzspieler, der den anderen ersetzt).
\end{eselsbruecke}

\section{Das ökonomische Prinzip}

Die Bedürfnisse der Menschen sind \emph{unbegrenzt}, die Güter dagegen \emph{knapp}. Aus
diesem \textbf{Spannungsverhältnis} folgt: Wer wirtschaftet --- private Haushalte wie
Unternehmen --- muss planvoll mit knappen Gütern umgehen. Das \textbf{ökonomische Prinzip}
(Wirtschaftlichkeitsprinzip) kennt zwei Ausprägungen:

\begin{itemize}
  \item \textbf{Maximalprinzip} --- mit \emph{gegebenen Mitteln} soll ein \emph{maximaler
    Erfolg} erzielt werden.
  \item \textbf{Minimalprinzip} --- ein \emph{gegebenes Ziel} soll mit \emph{minimalem
    Aufwand} erreicht werden.
\end{itemize}

\begin{beispiel}
\emph{Maximalprinzip:} Aus einer festen Menge Holz (gegebene Mittel) sollen möglichst
\emph{viele} Schreibtische (maximaler Erfolg) produziert werden. \\
\emph{Minimalprinzip:} Eine bestellte Menge Schreibtische (gegebenes Ziel) soll mit
möglichst \emph{wenig} Holz (minimaler Aufwand) hergestellt werden.
\end{beispiel}

\begin{merksatz}
Beide Prinzipien lassen sich \emph{nicht} gleichzeitig anwenden. Immer ist entweder das
Mittel fest vorgegeben (Maximalprinzip) \emph{oder} das Ziel (Minimalprinzip) --- nie
beides zugleich.
\end{merksatz}

\section{Produktionsfaktoren}

Güter müssen meist erst produziert werden. Alle an der Produktion beteiligten Menschen und
Güter heißen \textbf{Produktionsfaktoren}.

\subsection{Volkswirtschaftliche Produktionsfaktoren}

\begin{itemize}
  \item \textbf{Boden} --- die zu wirtschaftlichen Zwecken genutzte Natur:
    \emph{Anbauboden} (Land-/Forstwirtschaft), \emph{Abbauboden} (Rohstoffe, Bodenschätze),
    \emph{Standortboden} (Betriebsstätte).
  \item \textbf{Arbeit} --- jede geistige und körperliche Tätigkeit, mit der Einkommen
    erzielt wird. (Mäht der Hausmeister zu Hause den Rasen, zählt das \emph{nicht} dazu.)
  \item \textbf{Kapital} --- entsteht durch \emph{Sparen} (Konsumverzicht): \emph{Geldkapital}
    wird investiert und wird so zu \emph{Sachkapital} (z.\,B.\ Maschinen). Kapital gilt als
    \emph{derivativer} (abgeleiteter) Faktor --- es entsteht erst aus Boden und Arbeit.
\end{itemize}

\subsection{Betriebswirtschaftliche Produktionsfaktoren}

Im Betrieb unterscheidet man \textbf{Elementarfaktoren} und \textbf{dispositive Faktoren}:

\begin{center}
\begin{tabular}{p{6.7cm} p{6.7cm}}
\toprule
\textbf{Elementarfaktoren} & \textbf{Dispositive Faktoren}\\
\midrule
ausführende Arbeitskräfte & Leitung\\
Betriebsmittel (Maschinen, Fahrzeuge, Werkzeuge) & Planung\\
Werkstoffe (Roh-, Hilfs-, Betriebsstoffe, Fremdbauteile) & Organisation\\
 & Überwachung (Kontrolle)\\
\bottomrule
\end{tabular}
\end{center}

Die \emph{Werkstoffe} gliedern sich weiter: \textbf{Rohstoffe} sind Hauptbestandteil des
Produkts (Holz im Schreibtisch), \textbf{Hilfsstoffe} sind Nebenbestandteile (Leim,
Schrauben), \textbf{Betriebsstoffe} gehen nicht in das Produkt ein (Strom, Benzin, Öl).

Ziel ist die \textbf{Minimalkostenkombination}: die betriebswirtschaftlichen
Produktionsfaktoren so zu kombinieren, dass das Ergebnis mit minimalen Kosten erreicht
wird. Lassen sich Faktoren gegeneinander austauschen (z.\,B.\ Arbeit durch Maschine), sind
sie \emph{substitutional}; ist die Kombination technisch festgelegt, sind sie
\emph{limitational}.

\section{Der Wirtschaftskreislauf}

Der \textbf{Wirtschaftskreislauf} ist ein Modell, das die volkswirtschaftlichen
Zusammenhänge als Kreislauf darstellt. Er zeigt: Zwischen den Wirtschaftssektoren fließen
ständig Güter und Geld --- in \emph{entgegengesetzter} Richtung.

\subsection{Der einfache Wirtschaftskreislauf}

Betrachtet werden nur zwei Sektoren: die \emph{privaten Haushalte} und die
\emph{Unternehmen}. Annahmen (Prämissen): Nur die Haushalte besitzen die
Produktionsfaktoren; die Unternehmen entgelten deren Nutzung (Lohn/Gehalt für Arbeit,
Miete/Pacht für Boden, Zinsen für Kapital). Die Haushalte geben ihr gesamtes Einkommen für
Konsumgüter aus, die die Unternehmen herstellen.

\begin{itemize}
  \item \textbf{Güterstrom:} Produktionsfaktoren fließen von den Haushalten zu den
    Unternehmen; Konsumgüter und Dienstleistungen fließen zurück zu den Haushalten.
  \item \textbf{Geldstrom:} Einkommen/Entlohnung fließt zu den Haushalten; deren Ausgaben
    für Konsum fließen zu den Unternehmen.
\end{itemize}

Geld- und Güterstrom \emph{gleichen sich aus} (geschlossener Kreislauf). Da keine
Investitionen enthalten sind, nennt man ihn auch \emph{stationären} Wirtschaftskreislauf.

\subsection{Der erweiterte Wirtschaftskreislauf}

Nun sparen die Haushalte einen Teil ihres Einkommens. Das Modell wird um den Sektor
\textbf{Banken} erweitert: Über sie fließen Ersparnisse als \emph{Kredite} zu den
Unternehmen, die damit \emph{investieren} (z.\,B.\ in neue Maschinen). Im Modell ist das
Sparen der Haushalte genauso groß wie die Investitionen der Unternehmen. Weil so
Weiterentwicklung möglich wird, heißt er auch \emph{evolutionärer} Wirtschaftskreislauf.

\subsection{Der vollständige und der offene Wirtschaftskreislauf}

Kommt der \textbf{Staat} hinzu, entsteht der \emph{vollständige} Wirtschaftskreislauf. Der
Staat bezieht \emph{Steuern} von Haushalten und Unternehmen, zahlt \emph{Löhne/Gehälter}
und \emph{Sozialleistungen}, \emph{subventioniert} Unternehmen und konsumiert selbst Güter
und Dienstleistungen.

Bezieht man schließlich das \textbf{Ausland} mit ein, entsteht der Wirtschaftskreislauf
einer \emph{offenen} Volkswirtschaft: Güter, Dienstleistungen und Zahlungen werden
\emph{importiert} und \emph{exportiert}.

\begin{merksatz}
Vier Ausbaustufen: \textbf{einfach} (Haushalte $\leftrightarrow$ Unternehmen) $\rightarrow$
\textbf{erweitert} (+ Banken) $\rightarrow$ \textbf{vollständig} (+ Staat) $\rightarrow$
\textbf{offen} (+ Ausland). Geldstrom und Güterstrom laufen stets gegenläufig.
\end{merksatz}

\section{Zusammenfassung}

\begin{enumerate}
  \item \textbf{Bedürfnisse} sind Mangelgefühle; nach Dringlichkeit: Existenz-, Kultur-,
    Luxusbedürfnisse; ergänzt um materiell/immateriell, akut/latent, individual/kollektiv
    und die Pyramide nach Maslow.
  \item Über den \textbf{Bedarf} (mit Kaufkraft) wird ein Bedürfnis zur \textbf{Nachfrage};
    auf dem \textbf{Markt} bildet sich aus Angebot und Nachfrage der \textbf{Preis}.
  \item \textbf{Güter}: freie vs.\ knappe; Sachgüter/Dienstleistungen/Rechte;
    Konsum-/Produktions-; Gebrauchs-/Verbrauchs-; Komplementär-/Substitutionsgüter.
  \item Das \textbf{ökonomische Prinzip} kennt Maximal- (gegebene Mittel) und Minimalprinzip
    (gegebenes Ziel).
  \item \textbf{Produktionsfaktoren}: volkswirtschaftlich Boden, Arbeit, Kapital;
    betriebswirtschaftlich Elementar- (Arbeit, Betriebsmittel, Werkstoffe) und dispositive
    Faktoren.
  \item Der \textbf{Wirtschaftskreislauf} wächst von einfach (Haushalte/Unternehmen) über
    erweitert (Banken) und vollständig (Staat) bis offen (Ausland).
\end{enumerate}

\newpage

\section*{Glossar}
\addcontentsline{toc}{section}{Glossar}

\glossareintrag{Bedürfnis}{Gefühl eines Mangels mit dem Wunsch, ihn zu beseitigen --- Ausgangspunkt allen Wirtschaftens.}

\glossareintrag{Existenzbedürfnis}{Grundbedürfnis zur Sicherung des Überlebens (Essen, Trinken, Wohnung).}

\glossareintrag{Kulturbedürfnis}{Vom Kulturkreis geprägtes, als üblich angesehenes Bedürfnis (Bildung, Kino, Kleidung).}

\glossareintrag{Luxusbedürfnis}{Nicht notwendiger Anspruch, den sich nicht jeder leisten kann (Yacht, Villa).}

\glossareintrag{Individualbedürfnis}{Bedürfnis, das von einem einzelnen Menschen ausgeht.}

\glossareintrag{Kollektivbedürfnis}{Bedürfnis, das aus dem Zusammenleben entsteht (Umweltschutz, Autobahnnetz).}

\glossareintrag{Bedürfnispyramide (Maslow)}{Stufenmodell der Bedürfnisse von den physiologischen Grundbedürfnissen bis zur Selbstverwirklichung.}

\glossareintrag{Bedarf}{Mit Kaufkraft (Einkommen) ausgestattetes Bedürfnis.}

\glossareintrag{Nachfrage}{Marktwirksames Bedürfnis; der tatsächliche Kauf eines Gutes.}

\glossareintrag{Markt}{Ort des Zusammentreffens von Angebot und Nachfrage; hier bildet sich der Preis.}

\glossareintrag{Freie Güter}{In großer Menge vorhandene Güter ohne Preis (Luft, Sonne, Meerwasser).}

\glossareintrag{Knappe Güter}{Wirtschaftsgüter: begrenzt verfügbar, verursachen Kosten, haben einen Preis.}

\glossareintrag{Dienstleistung}{Immaterielles Gut; Handlung, durch die ein Nutzen entsteht (Reparatur, Beratung).}

\glossareintrag{Konsumgut}{Gut, das von Endverbrauchern (Konsumenten) genutzt wird.}

\glossareintrag{Produktionsgut}{Gut, das von Unternehmen im Herstellungsprozess eingesetzt wird.}

\glossareintrag{Komplementärgüter}{Güter, die sich ergänzen und gemeinsam nachgefragt werden (Maschine \& Strom).}

\glossareintrag{Substitutionsgüter}{Güter, die sich gegenseitig ersetzen (Butter oder Margarine).}

\glossareintrag{Ökonomisches Prinzip}{Grundsatz wirtschaftlichen Handelns; als Maximal- oder Minimalprinzip.}

\glossareintrag{Produktionsfaktoren}{An der Produktion beteiligte Menschen und Güter; volkswirtschaftlich Boden, Arbeit, Kapital.}

\glossareintrag{Wirtschaftskreislauf}{Modell der volkswirtschaftlichen Güter- und Geldströme zwischen den Sektoren.}

\end{document}
